\documentclass[10pt, aspectratio=169]{beamer}

% --- TEMA Y ESTILO ---
\usetheme{Madrid}
\usecolortheme{whale}

% --- PAQUETES ---
\usepackage[utf8]{inputenc}
\usepackage[spanish, es-tabla]{babel}
\usepackage{amsmath, amssymb}
\usepackage{siunitx}
\usepackage{fontawesome5}
\usepackage{booktabs}

\sisetup{
  output-decimal-marker = {,},
  per-mode = symbol
}

% --- INFORMACIÓN DEL CURSO ---
\title[Circuitos III - Dinámica Tribal]{Física de Semiconductores y Unión P-N}
\subtitle{Sesión 2: Aprendizaje Cooperativo (Método Jigsaw)}
\author{M.Sc. Ing. Bernardo Quiroga Turdera}
\institute[UCB Tarija]{Universidad Católica Boliviana ``San Pablo'' -- Sede Tarija}
\date{11 de agosto de 2026}

\begin{document}

\begin{frame}
  \titlepage
\end{frame}

\begin{frame}{Dinámica de Hoy: El Desafío Tribal}
  \begin{block}{Reglas del Juego}
    Hoy no habrá 90 minutos de cátedra magistral. El conocimiento lo construirán ustedes.
  \end{block}
  
  \vspace{0.3cm}
  \textbf{Fases de la Misión:}
  \begin{enumerate}
      \item \textbf{Enclave Tribal (25 min):} Su Tribu investiga y domina un concepto de la física de semiconductores.
      \item \textbf{Cruzada de Embajadores (30 min):} El Líder se queda a enseñar. Los Embajadores rotan de mesa en mesa para aprender y enseñar al resto.
      \item \textbf{Síntesis y Blindaje (15 min):} Ajuste matemático de ecuaciones y corrección de mitos (Docente).
      \item \textbf{Exit Ticket (10 min):} Preguntas al azar. La tribu depende del miembro menos preparado para ganar puntos PFI.
  \end{enumerate}
\end{frame}

\begin{frame}{Tribu 1: El Origen de la Carga}
  \begin{block}{\faMicrochip \quad Misión 1: Dopaje y Transporte}
    \textbf{Objetivo:} Explicar cómo transformamos la arena (Silicio) en conductor controlado.
    \vspace{0.2cm}
    \begin{itemize}
        \item Diferencia física real entre Tipo N y Tipo P (Ley de acción de masas: $n \cdot p = n_i^2$).
        \item Diferencia entre Corriente de \textbf{Deriva} (Drift) y de \textbf{Difusión}.
        \item Explicar la ecuación: $J_{\text{drift}} = q(n\mu_n + p\mu_p)E$.
    \end{itemize}
  \end{block}
  \vspace{0.3cm}
  \textbf{\faChalkboard \quad Reto en Pizarra:} Dibujar la red cristalina con un átomo de Fósforo (Tipo N) y uno de Boro (Tipo P) con sus portadores.
\end{frame}

\begin{frame}{Tribu 2: La Frontera Invisible}
  \begin{block}{\faProjectDiagram \quad Misión 2: Unión P-N y el Potencial $V_0$}
    \textbf{Objetivo:} Explicar qué ocurre exactamente al unir Tipo P con Tipo N.
    \vspace{0.2cm}
    \begin{itemize}
        \item Formación de la \textbf{Zona de Deplexión}. ¿Por qué está ``despoblada''?
        \item Surgimiento del Campo Eléctrico Interno ($E_0$).
        \item Explicar la ecuación del Potencial Integrado a temperatura ambiente:
        \begin{equation*}
            V_0 = V_T \ln \left( \frac{N_A N_D}{n_i^2} \right)
        \end{equation*}
    \end{itemize}
  \end{block}
  \vspace{0.1cm}
  \textbf{\faChalkboard \quad Reto en Pizarra:} Dibujar la unión P-N indicando los iones fijos y el vector $E_0$.
\end{frame}

\begin{frame}{Tribu 3: La Barrera de Potencial}
  \begin{block}{\faChartLine \quad Misión 3: La Ecuación de Shockley}
    \textbf{Objetivo:} La matemática de la curva $I$-$V$.
    \vspace{0.2cm}
    \begin{itemize}
        \item Zona de deplexión en \textbf{Polarización Inversa} (y qué es $I_S$).
        \item Zona de deplexión en \textbf{Polarización Directa}.
        \item Explicar término a término la Ecuación del Diodo:
        \begin{equation*}
            I_D = I_S \left( \exp\left(\frac{V_D}{\eta V_T}\right) - 1 \right)
        \end{equation*}
    \end{itemize}
  \end{block}
  \vspace{0.1cm}
  \textbf{\faChalkboard \quad Reto en Pizarra:} Dibujar la curva no lineal $I$-$V$ del diodo marcando la rodilla (\qty{0.7}{\volt}) y $I_S$.
\end{frame}

\begin{frame}{Tribu 4: Resistencia Dinámica}
  \begin{block}{\faWaveSquare \quad Misión 4: Modelo de Pequeña Señal ($r_d$)}
    \textbf{Objetivo:} Un diodo no es un interruptor ideal; es una resistencia variable.
    \vspace{0.2cm}
    \begin{itemize}
        \item ¿Qué es el Punto de Operación en DC ($Q$)?
        \item Demostrar analíticamente cómo se obtiene $r_d$:
        \begin{equation*}
            r_d = \frac{d V_D}{d I_D} \approx \frac{\eta V_T}{I_D}
        \end{equation*}
        \item Ejercicio: Si $I_D = \qty{2}{\milli\ampere}$ ($V_T = \qty{25.85}{\milli\volt}$, $\eta=1$), ¿cuánto vale $r_d$? ¿Y si $I_D = \qty{20}{\milli\ampere}$?
    \end{itemize}
  \end{block}
  \vspace{0.1cm}
  \textbf{\faChalkboard \quad Reto en Pizarra:} Dibujar la recta tangente en el punto Q.
\end{frame}

\begin{frame}{Síntesis y Próximos Pasos}
  \begin{block}{Conclusión del Docente}
      El comportamiento no lineal del diodo está gobernado físicamente por la ecuación de Shockley. Su resistencia en AC ($r_d$) es inversamente proporcional a la corriente DC.
  \end{block}
  
  \vspace{0.5cm}
  \textbf{Conexión con el PFI (Hito 1):}
  \begin{itemize}
      \item Los diodos Zener/Flyback cambian su resistencia dinámica ($r_d$) según la intensidad del pico inductivo.
      \item \textbf{Actividad asíncrona:} Correr el script en Python para comprobar cómo $r_d$ e $I_S$ varían con la temperatura en Tarija.
  \end{itemize}
\end{frame}

\end{document}